\chapter*{考研三卷种阅读路线}
\addcontentsline{toc}{chapter}{考研三卷种阅读路线}
本书同时保存基础计算、综合证明与提高材料。考研复习按卷种选择内容，不把全书目录视为每位考生的必修清单。本页按教育部教育考试院《2027年全国硕士研究生招生考试数学考试大纲》核读，复习其他年份时另对当年大纲。

共同主线包括极限连续、一元微积分、相应微分方程、多元微分与二重积分；各卷对细目的要求仍有区别。
\begin{center}
\begin{tabularx}{\textwidth}{@{}p{.32\textwidth}XXX@{}}
\toprule 入口&数学一&数学二&数学三\\\midrule
第5章空间几何&按大纲读&不作单独主线&不作单独主线\\
第7章三重积分、第8章&按大纲读&跳过&跳过\\
第9章数项与幂级数&按大纲读&跳过&按大纲读\\
第9章 Fourier 内容&按大纲读&跳过&跳过\\
经济应用与差分&按本卷范围选&按本卷范围选&对照本卷要求自查\\
\bottomrule
\end{tabularx}
\end{center}

“按大纲读”不等于整章所有提高题都在考试范围内。实数完备性证明体系、椭圆或复变方法、相空间与稳定性、一致收敛的深入理论等不作为统考复习主线；夹在这些专题里的初等计算、极限与判敛方法仍应按具体题目保留。数三的边际、弹性与差分应用见第\ref{r2:economics}、\ref{r2:difference}节；初等模型用于练习列式与求解，不把所有稳定性推广都视为统考必考内容。

\section*{从题目入口回查正文}
\begin{itemize}
\item \textbf{计算基础：}先读第1章的极限方法与第3章的不定积分选法；遇到条件不明再回查定义、余项和定理适用条件，避免一次堆入多种高级方法。
\item \textbf{中值定理：}由第\ref{sec:2:3}节回查辅助函数构造。先写连续、可导的区间条件，再核对端点关系与结论中的点所在区间。
\item \textbf{闭区域最值：}见第\ref{sec:6:4}—\ref{sec:6:5}节。依次列出内部驻点、各段边界的候选点、边界端点或角点，并统一比较函数值；约束不正则的点不能直接被乘子法略去。
\item \textbf{投影与积分限：}见第\ref{sec:5:5}、\ref{sec:7:1}—\ref{sec:7:3}及\ref{sec:8:3}节。先画区域，区分投影曲线、投影柱面和投影区域，消元后回查原约束与分支。
\item \textbf{积分应用：}第3章既有应用题可用来练习微元、密度、力与距离的对应关系。是否属于本卷训练任务，先对照大纲；不把所有物理应用都笼统写成三卷共同的高频题。
\item \textbf{Fourier（数学一）：}见第\ref{sec:9:5}节。先判断延拓与周期，再求系数，最后用周期延拓后的左右极限确定间断点及端点的和函数值。
\end{itemize}

复习是否完成，应以相应卷种的限时题目和错因复盘判断，不用“每年必考”“隔年必出”等未经逐年统计的说法替代范围核对。卷面内容比例约为：数学一、三的微积分部分60\%，数学二80\%；这些是卷面构成，不是章节出题概率。
官方版本入口：\href{https://xuanshu.hep.com.cn/front/book/findBookDetails?bookId=6aa1910be119ac97297a608f}{高等教育出版社：2027数学考试大纲}。本次核读原书第5页及相应卷种的微积分范围；核验记录另见随书修订说明。
